\section*{Capítulo 3}

\solucion{3.1} Rendimientos: $\ln(84/80)=4.879\,\%$; $\ln(78/84)=-7.411\,\%$; $\ln(81/78)=3.774\,\%$; $\ln(86/81)=5.990\,\%$. Media: $1.808\,\%$. Desviaciones respecto de la media: $3.071$, $-9.219$, $1.966$, $4.182$ (en \%). Suma de cuadrados: $9.43+84.99+3.87+17.49=115.78$. Varianza muestral: $115.78/3=38.59$; desviación estándar: $6.21\,\%$ mensual. Anualizada: $6.21\times\sqrt{12}=21.5\,\%$.

\solucion{3.2} En dólares: $2{,}000/3.50=571.43$ millones. TC $-5\,\%$: $571.43\times(-0.05/0.95)=-30.08$. TC $+5\,\%$: $571.43\times(0.05/1.05)=+27.21$. La aproximación lineal da $\pm28.57$: subestima la pérdida y sobrestima la ganancia.

\solucion{3.3} El valor en dólares de un sol es $1/S$. Cuando $S$ baja, $1/S$ sube cada vez más rápido (la curva $1/S$ es convexa): un mismo porcentaje de caída del dólar hace subir el valor del sol más de lo que lo baja una subida del mismo porcentaje. Para quien debe soles, eso se traduce en pérdidas mayores que las ganancias simétricas.

\solucion{3.4} Valores presentes al 8\,\%: $6/1.08=5.556$; $6/1.08^2=5.144$; $106/1.08^3=84.146$. Precio: $94.846$. $D_{Mac}=(5.556+2\times5.144+3\times84.146)/94.846=268.28/94.846=2.829$ años. $D_{mod}=2.829/1.08=2.619$. $\Delta P\approx-2.619\times0.005\times94.846=-1.24$; el precio pasa a unos 93.60.

\solucion{3.5} $8.46\,\%\times\sqrt{12}=29.3\,\%$. Bajo normalidad, el 68\,\% de los rendimientos log anuales estaría entre $-29.3\,\%$ y $+29.3\,\%$ (suponiendo media cercana a cero).

\solucion{3.6} No. La deuda a tasa fija elimina el riesgo de que suban los \emph{intereses que se pagan}, pero el valor razonable de esa deuda cambia con las tasas (US\$ 166 millones por 100 pb en los bonos) y, sobre todo, la empresa tendrá que \emph{refinanciar} deudas que vencen (el CESCE, reclasificado al corto plazo, y las líneas de corto plazo) a las tasas que rijan entonces. Ese es su riesgo de tasa: un riesgo de refinanciación.

\solucion{3.7} $z\sigma\sqrt T=1.645\times0.06\times0.5=0.04935$. TC adverso: $3.40\,e^{-0.04935}=3.2363$. VaR $=500/3.2363-500/3.40=154.50-147.06=7.44$ millones de dólares.

\solucion{3.8} $\sqrt{80^2+50^2+2\times0.35\times80\times50}=\sqrt{6{,}400+2{,}500+2{,}800}=\sqrt{11{,}700}=108.2$ millones, frente a una suma simple de 130. La diferencia (21.8) es el beneficio de diversificación: es poco probable que ambos riesgos se materialicen con toda su fuerza a la vez.
